%================================================================
\chapter{Evaluación y Registro de Progreso}
\label{ch:assessment}
%================================================================

\begin{Form}

Este capítulo es un formulario PDF rellenable dirigido al profesorado
y al personal que se incorpora como colaborador del repositorio de
currículo del Departamento de Matemáticas. Cada sección corresponde
a una categoría de preparación: Obligatorio, Esencial, Deseable y
Opcional. Las casillas de verificación registran la finalización
binaria de los pasos de verificación; los campos de texto registran
respuestas escritas breves que demuestran comprensión conceptual.
Utilice este capítulo como una lista de verificación activa durante
su incorporación.

\vs
\noindent
\textbf{Guardar sus respuestas requiere Adobe Acrobat Reader}
(gratuito para Windows, macOS y Linux; disponible en
\url{https://get.adobe.com/reader/}). Otros visores se comportan
de la siguiente manera:

\begin{itemize}
    \item \textbf{Vista Previa de macOS:} muestra los campos del
          formulario y permite escribir en ellos, pero \emph{no guarda
          los datos ingresados}. Si cierra y vuelve a abrir el archivo
          en Vista Previa, sus respuestas se perderán.
    \item \textbf{Google Chrome (visor de PDF integrado):} muestra y
          acepta entradas en los campos del formulario, pero \emph{no
          guarda los datos}. Imprimir a PDF desde Chrome tampoco
          preservará el contenido de los campos.
    \item \textbf{La mayoría de los visores de PDF de Linux} (Evince,
          Okular, Zathura): \emph{ignoran por completo los campos del
          formulario}. Los campos no aparecerán interactivos.
    \item \textbf{Adobe Acrobat Reader:} guarda, imprime y exporta
          completamente los datos del formulario. Use \textbf{Archivo
          $>$ Guardar} (no Guardar como) después de completar cada
          sección.
\end{itemize}

\vs
\noindent
\textbf{Cómo interpretar las cuatro categorías:}

\begin{description}
    \item[Obligatorio] Tareas básicas del repositorio. Estas deben
          completarse antes de que un colaborador pueda realizar
          cualquier cambio en el documento de currículo.
    \item[Esencial] Configuración del entorno local: \LaTeX{}, Python,
          Git y Visual Studio Code. Estos deben funcionar
          correctamente antes de comenzar cualquier trabajo de autoría
          o de programación.
    \item[Deseable] Comprensión conceptual de las herramientas
          introducidas en los capítulos de infraestructura. Completar
          estos elementos demuestra estar preparado para trabajar de
          manera independiente en el repositorio.
    \item[Opcional] Infraestructura local de modelos grandes de
          lenguaje (Ollama). No es obligatorio para la autoría del
          currículo; es relevante para los colaboradores que deseen
          usar modelos de lenguaje locales en su flujo de trabajo.
\end{description}
\vs

%================================================================
\section{Obligatorio: Incorporación al Repositorio}
\label{sec:assess-required}
%================================================================

Los siguientes elementos verifican que usted puede recuperar,
compilar y contribuir al repositorio de currículo. Todos los
elementos deben completarse antes de editar cualquier archivo fuente.

\vs

\noindent\CheckBox[name=r1,width=0.4cm,height=0.4cm]{}~\textbf{R1.}
El \texttt{git clone} del repositorio se completa sin errores y la
copia de trabajo local contiene la carpeta \texttt{curriculum/}.
(Véase el Capítulo~\ref{ch:git}, Sección~\ref{subsec:multiple-remotes}.)

\vs

\noindent\CheckBox[name=r2,width=0.4cm,height=0.4cm]{}~\textbf{R2.}
Ejecutar \texttt{pdflatex MathCurriculum.tex} desde dentro de la
carpeta \texttt{curriculum/} produce \texttt{MathCurriculum.pdf}
sin errores fatales.
(Véase el Capítulo~\ref{ch:latex}, Sección~\ref{sec:first-latex-doc}.)

\vs

\noindent\CheckBox[name=r3,width=0.4cm,height=0.4cm]{}~\textbf{R3.}
Ejecutar \texttt{gh} (o \texttt{gh.bat} en Windows) con un mensaje
de confirmación prepara, confirma y envía exitosamente un cambio al
repositorio, y \texttt{VERSION.txt} se actualiza con el nuevo hash.
(Véase el Capítulo~\ref{ch:git}.)

\vs

\noindent\CheckBox[name=r4,width=0.4cm,height=0.4cm]{}~\textbf{R4.}
El hash de confirmación en \texttt{VERSION.txt} coincide con el
valor devuelto por \texttt{git rev-parse --short HEAD} en su
repositorio local.
(Véase el Capítulo~\ref{ch:git}.)

\vs

\noindent\textbf{R5.} Describa la estructura de dos partes del
documento de currículo. ¿Qué está contenido en los capítulos de
infraestructura y qué está contenido en los capítulos de unidades
del curso? ¿Cómo se identifican los capítulos de unidades en el
árbol de código fuente?
(Véase la Introducción.)\\[4pt]
\TextField[name=r5,width=\linewidth,height=3cm,multiline=true,
           bordercolor=black]{}

%================================================================
\section{Esencial: Configuración del Entorno y Cadena de Herramientas}
\label{sec:assess-essential}
%================================================================

Marque cada casilla cuando el paso de verificación correspondiente
tenga éxito. Todos los elementos de esta sección deben estar marcados
antes de comenzar cualquier trabajo de autoría o de programación en
el repositorio.

\vs

\noindent\CheckBox[name=e1,width=0.4cm,height=0.4cm]{}~\textbf{E1.}
\texttt{pdflatex --version} devuelve un número de versión en la
terminal.
(Véase el Capítulo~\ref{ch:environment-setup}.)

\vs

\noindent\CheckBox[name=e2,width=0.4cm,height=0.4cm]{}~\textbf{E2.}
\texttt{python --version} devuelve un número de versión en la terminal.
(Véase el Capítulo~\ref{ch:environment-setup}, Sección~\ref{sec:python}.)

\vs

\noindent\CheckBox[name=e3,width=0.4cm,height=0.4cm]{}~\textbf{E3.}
\texttt{git --version} devuelve un número de versión.
(Véase el Capítulo~\ref{ch:environment-setup}, Sección~\ref{sec:git}.)

\vs

\noindent\CheckBox[name=e4,width=0.4cm,height=0.4cm]{}~\textbf{E4.}
\texttt{code .} abre Visual Studio Code desde una ventana de terminal.
(Véase el Capítulo~\ref{ch:environment-setup}, Sección~\ref{sec:vscstart}.)

\vs

\noindent\CheckBox[name=e5,width=0.4cm,height=0.4cm]{}~\textbf{E5.}
El entorno virtual de Python se activa sin errores, y los paquetes
listados en \texttt{pyproject.toml} (o \texttt{requirements.txt})
se importan correctamente en ese entorno.
(Véase el Capítulo~\ref{ch:python-environments},
Sección~\ref{sec:project-libraries}.)

\vs

\noindent\CheckBox[name=e6,width=0.4cm,height=0.4cm]{}~\textbf{E6.}
Ejecutar \texttt{latexmk -pdf MathCurriculum.tex} produce un PDF
correctamente paginado con una tabla de contenido y todas las
referencias cruzadas resueltas.
(Véase el Capítulo~\ref{ch:latex}.)

%================================================================
\section{Deseable: Comprensión Técnica}
\label{sec:assess-desirable}
%================================================================

Estos elementos demuestran una comprensión funcional de las
herramientas introducidas en los capítulos de infraestructura. Las
respuestas escritas no necesitan ser extensas; dos o tres oraciones
cada una son suficientes.

\vs

\noindent\textbf{D1.} ¿Cuál es el directorio de trabajo actual cuando
usted abre por primera vez una terminal? Escriba el comando que
utilizaría para navegar hasta la carpeta \texttt{curriculum/} dentro
del repositorio clonado.
(Véase el Capítulo~\ref{ch:cli-scripting},
Sección~\ref{sec:what-is-a-file}.)\\[4pt]
\TextField[name=d1,width=\linewidth,height=2.5cm,multiline=true,
           bordercolor=black]{}

\vs

\noindent\textbf{D2.} ¿Cuál es la diferencia entre un entorno virtual
de Python y la instalación global de Python? Proporcione un escenario
en el contexto de este repositorio donde necesitaría un entorno
virtual separado.
(Véase el Capítulo~\ref{ch:python-environments}.)\\[4pt]
\TextField[name=d2,width=\linewidth,height=2.5cm,multiline=true,
           bordercolor=black]{}

\vs

\noindent\textbf{D3.} ¿Por qué \texttt{git push} requiere un
\texttt{git commit} previo? ¿Cuál es la diferencia conceptual entre
las dos operaciones, y por qué el script \texttt{gh.bat} realiza
ambas en secuencia?
(Véase el Capítulo~\ref{ch:git}.)\\[4pt]
\TextField[name=d3,width=\linewidth,height=2.5cm,multiline=true,
           bordercolor=black]{}

\vs

\noindent\textbf{D4.} En \LaTeX{}, ¿cuál es la diferencia práctica
entre \verb|\[...\]| y \verb|$$...$$|? ¿Qué forma debería usarse en
los nuevos archivos fuente del currículo y por qué?
(Véase el Capítulo~\ref{ch:latex}, Sección~\ref{sec:latex-math}.)\\[4pt]
\TextField[name=d4,width=\linewidth,height=2.5cm,multiline=true,
           bordercolor=black]{}

\vs

\noindent\textbf{D5.} ¿Qué hace el script \texttt{regen.py} cuando se
ejecuta para una unidad de curso determinada? Describa la entrada
que lee, la salida que produce y el papel del diccionario
\texttt{ASSESSMENTS} en el control de la salida.
(Véase la carpeta \texttt{Script/} del repositorio y la documentación
asociada.)\\[4pt]
\TextField[name=d5,width=\linewidth,height=2.5cm,multiline=true,
           bordercolor=black]{}

\vs

\noindent\textbf{D6.} Explique el esquema de etiquetas utilizado para
las referencias cruzadas en los archivos fuente del currículo. ¿Cuál
es el formato de una etiqueta de sección, cómo se deriva el mnemónico
de cuatro caracteres y cómo se calcula el sufijo hash?
(Véase la documentación de la canalización del currículo.)\\[4pt]
\TextField[name=d6,width=\linewidth,height=2.5cm,multiline=true,
           bordercolor=black]{}

%================================================================
\section{Opcional: Infraestructura Local de LLM (Ollama)}
\label{sec:assess-optional}
%================================================================

Estos elementos se aplican solo a los colaboradores que instalan y
configuran la infraestructura local del modelo grande de lenguaje
Ollama. No son obligatorios para la autoría del currículo ni para el
mantenimiento del repositorio.

\vs

\noindent\CheckBox[name=o1,width=0.4cm,height=0.4cm]{}~\textbf{O1.}
Ollama está instalado y \texttt{ollama list} devuelve al menos un
modelo en la terminal.
(Véase el Capítulo~\ref{ch:ollama-infrastructure},
Sección~\ref{sec:install}.)

\vs

\noindent\CheckBox[name=o2,width=0.4cm,height=0.4cm]{}~\textbf{O2.}
\texttt{connectivity\_test.py} aprueba todas las verificaciones
(protocolo de enlace TCP, lista de modelos, inferencia).
(Véase el Capítulo~\ref{ch:ollama-infrastructure}.)

\vs

\noindent\textbf{O3.} Explique el equilibrio entre el tamaño de la
ventana de contexto y el consumo de memoria en Ollama. ¿Por qué
\texttt{OLLAMA\_NUM\_CTX=8192} es un valor predeterminado seguro
para modelos grandes en hardware con VRAM limitada?
(Véase el Capítulo~\ref{ch:ollama-infrastructure},
Sección~\ref{sec:memory}.)\\[4pt]
\TextField[name=o3,width=\linewidth,height=3cm,multiline=true,
           bordercolor=black]{}

\end{Form}
